\answer
\begin{enumerate}

%%--- 章末問題 4.1 ---
\item[\ref{ex4.1}]
定理\ref{th4.1}のボルト秒平衡より
\begin{equation*}
DT\,V_0 + (T - DT)(V_0 - V_1) = 0
\end{equation*}
これを整理すると$TV_0 = (1-D)TV_1$，よって
\begin{equation*}
V_1 = \frac{V_0}{1-D}
\end{equation*}
$0 < D < 1$では$V_1 > V_0$，すなわち\textbf{出力電圧が入力より高くなる}。
オン期間に蓄えた磁気エネルギーを，オフ期間に入力電圧へ上乗せして送り出すためである。
これが昇圧コンバータ（5章）の原理である。

%%--- 章末問題 4.2 ---
\item[\ref{ex4.2}]
オン期間のインダクタ電圧は$v_L = V_0 - V_1 = (1-D)V_0$で一定だから，
式\ref{eq4.5}より電流はオン期間$DT$の間に
\begin{equation*}
\mathit{\Delta}i_L = \frac{1}{L}(1-D)V_0 \times DT = \frac{D(1-D)V_0 T}{L}
\end{equation*}
だけ直線的に増える。これが$I_{\max} - I_{\min}$である。
\rev{数値を代入すると，$T = 1/f = 10\,\mathrm{\mu s}$だから}
\begin{equation*}
\mathit{\Delta}i_L = \frac{0.5\times0.5\times12\times10\times10^{-6}}{100\times10^{-6}}
= 0.3\,\mathrm{A}
\end{equation*}
\rev{リプルは$L$に反比例し，$f$を上げても小さくなる。}
また$D(1-D)$は$D = 0.5$のとき最大になるから，この数値は最悪条件での値である。

\rev{$\mathit{\Delta}i_L \le 0.1$\,Aにするには，上の式を$L$について解いて}
\begin{equation*}
\rev{L \ge \frac{D(1-D)V_0 T}{\mathit{\Delta}i_L}
= \frac{0.5\times0.5\times12\times10\times10^{-6}}{0.1} = 300\,\mathrm{\mu H}}
\end{equation*}
\rev{すなわち$300\,\mathrm{\mu H}$以上にすればよい。
負荷電流が$2$\,Aなら$I_{\mathrm{avg}} = 2$\,Aだから，式\ref{eq4.dWL}より1周期に往復するエネルギーは}
\begin{equation*}
\rev{\mathit{\Delta}W_L = L\,I_{\mathrm{avg}}\,\mathit{\Delta}i_L
= 300\times10^{-6}\times2\times0.1 = 60\,\mathrm{\mu J}}
\end{equation*}
\rev{一方$L = 100\,\mathrm{\mu H}$では$\mathit{\Delta}i_L = 0.3$\,Aで
$\mathit{\Delta}W_L = 100\times10^{-6}\times2\times0.3 = 60\,\mathrm{\mu J}$ ― 同じである。
往復するエネルギーは，負荷へ送る電力$P = V_1I_{\mathrm{avg}} = 6\times2 = 12$\,Wのうち
オフ期間の分$(1-D)P/f = 0.5\times12/10^{5} = 60\,\mathrm{\mu J}$で決まり，$L$にはよらない（\ref{sec4.2}節）。
$L$を3倍にすると，同じ$60\,\mathrm{\mu J}$の往復を3倍の$LI_{\mathrm{avg}}$で受け持つので，
式\ref{eq4.dWL2}のとおりリプルだけが$1/3$に減る。
インダクタを大きくしても「運ぶエネルギー」は変わらず，変わるのは「電流の振れ幅」である。}

%%--- 章末問題 4.3 ---
\item[\ref{ex4.3}]
オン期間のキャパシタ電流は$i_C = I_0 - I_1 = (1-D)I_0$で一定だから，
式\ref{eq4.9}より電圧はオン期間$DT$の間に
\begin{equation*}
\mathit{\Delta}v_C = \frac{1}{C}(1-D)I_0 \times DT = \frac{D(1-D)I_0 T}{C}
\end{equation*}
だけ直線的に上がる。数値を代入して
\begin{equation*}
\mathit{\Delta}v_C = \frac{0.5\times0.5\times2\times10\times10^{-6}}{10\times10^{-6}}
= 0.5\,\mathrm{V}
\end{equation*}
章末問題\ref{ex4.2}の式で$V_0 \to I_0$，$L \to C$，
$\mathit{\Delta}i_L \to \mathit{\Delta}v_C$と読み替えたものにちょうど一致する。
インダクタの電流リプルとキャパシタの電圧リプルは双対（\TABref{tab4.1}）であり，
片方を導けばもう片方は読み替えだけで得られる。

%%--- 章末問題 4.4 ---
\item[\ref{ex4.4}]
(a) 式\ref{eq4.4}と式\ref{eq4.8}より
\begin{equation*}
W_L = \frac{1}{2}\times100\times10^{-6}\times10^2 = 5\,\mathrm{mJ}，\qquad
W_C = \frac{1}{2}\times100\times10^{-6}\times400^2 = 8\,\mathrm{J}
\end{equation*}
(b) 1周期は$T = 10\,\mathrm{\mu s}$だから，毎周期$5$\,mJを受け渡すと平均電力は
\begin{equation*}
\rev{P = W_L f = 5\times10^{-3}\times10^{5} = 500\,\mathrm{W}}
\end{equation*}
1回に蓄えられる量はわずか$5$\,mJでも，受け渡しを毎秒10万回繰り返せば
$500$\,Wもの電力を運べる。
\rev{ただしこれは，電流が毎周期$0$から$10$\,Aまで振れて蓄えを空にする運転
（電流が途切れる境界，または不連続の運転）での値である。
電流が途切れない運転では，往復するのは式\ref{eq4.dWL}の$\mathit{\Delta}W_L$だけで，
たとえば章末問題\ref{ex4.2}の条件（$I_{\mathrm{avg}} = 2$\,A，$\mathit{\Delta}i_L = 0.3$\,A）では
$60\,\mathrm{\mu J}$，電力にして$6$\,W分（負荷電力$12$\,Wのオフ期間の分）にすぎない。}
スイッチング周波数を高くするほど，同じインダクタ（あるいはもっと小さなインダクタ）で
大きな電力を扱える ― 電源装置の小型化が高周波化とともに進んできた理由である。

%%--- 章末問題 4.5 ---
\item[\ref{ex4.5}]
（シミュレーションの確認ポイント）
(a) 出力電圧の平均は$DV_0 = 6$\,Vに落ち着く。
起動直後は振動（過渡状態）が見えるが，やがて定常状態に収束する。
波形の平均を取る区間は定常状態に入った後に選ぶこと。
(b) リプル電圧はほぼ$C$に反比例するから，$C = 100\,\mathrm{\mu F}$では
約$1/10$になる。
(c) \rev{$f = 5$\,kHzは$f_0 = 1/(2\pi\sqrt{LC}) \approx 5.0$\,kHzとほぼ一致し，}
スイッチング成分が共振で増幅されて出力は大きく振動する。
\rev{「$f_0 \ll f$」というフィルタ設計の条件（式\ref{eq4.13}）が}
崩れるとどうなるかを目で確かめられる。

\end{enumerate}
